\section{Related Work}
This section reviews both traditional and deep learning methods for point cloud registration. The recent developments of NeRF on LiDAR applications and image/RGB-based NeRF for BA are also reviewed, followed by a focused review on volume sampling for NeRFs.

\subsection{Point Cloud Registration}
Point cloud registration has been one of the long-standing problems in LiDAR research, with applications including LiDAR odometry and 3D mapping \cite{Zhang2014LOAM}. Traditional techniques for point cloud registration, such as ICP~\cite{besl1992icp} or point to plane ICP~\cite{CHEN1992point2planeicp}, suffer from local minima, while deep learning-based methods~\cite{wang2023sghr,qin2023geotrans,lu2021hregnet}, although achieving state-of-the-art results, suffer from limited generalization and the need for a large amount of point cloud data. Meanwhile, some works study Bundle Adjustment (BA) of LiDAR 3D point clouds \cite{liu2022HBA,ruhnke2012highly} \cite{liu2021balm}. However, correcting 3D LiDAR points through classic BA methods is challenging due to the inherently unordered and sparse nature of point clouds. 

\subsection{LiDAR NeRF and Neural SDF}
Recent research in NeRF or radiance fields has extended from RGB images to LiDAR data, including novel LiDAR view synthesis~\cite{zheng2024lidar4D,jiang2025gslidar}, mapping and LiDAR odometry~\cite{deng2023nerfloam,isaacson2024lonerlidar}. While these works achieved impressive results, very few have discussed the capability of BA for NeRF. GeoNLF~\cite{xue2024geonlf} has explored the idea of LiDAR NeRF-BA by combining NeRF with a geometry optimizer for LiDAR sensor poses, however, it focused on designing the external architecture and overlooked the underlying mechanism for LiDAR NeRF-BA. LONER~\cite{isaacson2024lonerlidar} proposed a LiDAR NeRF-BA framework that jointly optimizes LiDAR poses and NeRF parameters. However, it did not discuss the influence of volume sampling on LiDAR NeRF-BA. 

A parallel line of research represents LiDAR scenes with neural signed distance fields (SDFs)~\cite{zhong2023shinemapping}\cite{pan2024pinslam}\cite{deng2023nerfloam}. While neural SDFs produce compact and accurate reconstructions, several of their standard design choices~\cite{takikawa2021neuralsdf} might not be suitable for LiDAR-BA. For example, phere tracing range rendering could lead to gradient explosion or vanishing during pose optimization due to recursive evaluation on the neural SDF. For these reasons, we build our LiDAR NeRF-BA on density-based volume rendering rather than implicit distance fields.

\subsection{RGB NeRF-BA}
While RGB NeRFs were originally designed for novel-view synthesis with known camera poses~\cite{mildenhall2020nerfOG}, their differentiable nature can also enable gradient-based BA to jointly refine camera poses~\cite{wang2022nerfmm,bian2023nopenerf}. Further research advancements have shown techniques, such as gradual activation of positional encoding frequencies~\cite{lin2021barf} and over-parameterization of camera poses~\cite{chen2023l2gnerf}, can significantly improve BA results. Similar ideas have been applied to image deblurring~\cite{wang2023badnerf}. Although these techniques could inspire LiDAR NeRF-BA, they do not consider the specific characteristics of LiDAR rays.

\subsection{Volume Sampling in NeRF}
Volume rendering of rays requires volume sampling in NeRF~\cite{mildenhall2020nerfOG}, which involves selecting a set of points along the ray to query the NeRF model for volume rendering. In the original NeRF implementation, hierarchical sampling is proposed to first render coarse ray samples before resampling around volumes with high ray termination probability. This allows for gradient-efficient NeRF model training. Later advancements in research have improved volume sampling to effectively sample conical frusta \cite{barron2021mipnerf,barron2023zipnerf} or cones on reflected surfaces \cite{verbin2024nerfcasting}. However, previous research on volume sampling focuses on improving image synthesis quality, with no existing research discussing its significance or design in LiDAR NeRF.

% Research on volume sampling focused on improving optimisation efficiency, anti-aliasing. But overlooked its significance in joint optimisation of LiDAR sensor pose.